%! The Victory of Orthogonality — Unit 7, Lesson 7.4 — Exit Ticket
\documentclass[10pt]{article}
\usepackage{linalg-article}
\usepackage{linalg-boxes}
\newcommand{\TallMath}[1]{$\displaystyle #1\rule[-1.4em]{0pt}{3.2em}$}
\newcommand{\uu}{\mathbf{u}}
\newcommand{\vv}{\mathbf{v}}
\newcommand{\xx}{\mathbf{x}}
\newcommand{\zero}{\mathbf{0}}
\newcommand{\T}{^{\top}}

\begin{document}

\pageheader{Unit 7, Lesson 7.4}{Exit Ticket}
\namedateperiod

\vspace{0.15in}

No notes. Show your reasoning. Consider the matrix
\[
  Q=\tfrac15\begin{bmatrix} 4 & -3\\ 3 & 4 \end{bmatrix}.
\]

\begin{enumerate}[leftmargin=1.8em, itemsep=15pt]
  \item \textbf{Is it orthogonal?} Compute $Q\T Q$. If $Q$ is orthogonal, write its inverse $Q^{-1}$ (no computation needed).
        \vspace{2.0cm}
  \item \textbf{Length.} A vector $\xx$ has $|\xx|=6$. Without finding $\xx$, what is $|Q\xx|$? Explain in one line why.
        \vspace{1.8cm}
  \item \textbf{Explain.} The SVD writes any matrix as $A=U\Sigma V\T$. (a) Why is this called ``rotate--stretch--rotate''? (b) Why does that make orthogonal matrices the ``victory'' of this unit? \writelines{2}
\end{enumerate}

\end{document}
